% --------------------------------------------------
% 7. CONCLUSION
% --------------------------------------------------
\section{Conclusion}
\label{sec:conclusion}

This paper has examined rider abandonment in Grand Tour cycling as an economic decision, using survival analysis to study 9,045 rider-race observations across the Tour de France, Giro d'Italia, and Vuelta a Espa\~{n}a from 2010 to 2025.

The Cox proportional hazard models provide three main findings. First, a rider's historical DNF rate is the strongest predictor of abandonment (HR $= 1.549$), consistent with the interpretation that past abandonment captures rider-specific expected continuation costs. The sub-sample analysis strengthens this: the effect is substantially larger for GC contenders (HR $= 2.280$) than for non-GC riders (HR $= 1.472$), suggesting heterogeneous cost structures. Second, riding in the Tour de France is associated with approximately 14\% lower abandonment hazard, consistent with the event-value hypothesis: higher prize money, media exposure, and UCI points raise the expected return from continued participation \parencite{blair2012sports}. This effect is significant for non-GC riders but not for GC contenders, suggesting the incentive gradient operates primarily through domestiques and support riders. Third, mountain stages raise daily abandonment hazard, consistent with a marginal-cost-of-effort interpretation.

These findings are consistent with a model in which abandonment is a rational decision shaped by the balance of incentives and costs, in the spirit of tournament theory \parencite{rosen1981superstars, lazear1981rank} and the labour market laboratory perspective \parencite{kahn2000labormarket}.

Several limitations should be acknowledged. The most important is the absence of direct financial data: rider salaries, team budgets, and prize money receipts are unavailable at the rider-race level. The economic interpretation rests on the assumption that the proxy variables are positively correlated with the underlying constructs. A second limitation is that the GC contender classification is a coarse binary indicator based on past results, which does not capture within-race changes in GC standing. Third, the dataset does not distinguish between types of abandonment: crashes, illness, and strategic withdrawal are coded identically.

Future work could extend the analysis in several directions. Incorporating prize-money or UCI-points data would allow a more direct test of the event-value hypothesis. Adding contract-year indicators would enable a test of career concerns theory \parencite{boeri2008italianjob}. A competing-risks model separating different abandonment causes would provide a more refined picture of the decision process. A counterfactual analysis estimating how GC standings would change if abandoned riders had continued (see Figure~\ref{fig:counterfactual} in the Appendix for a preliminary exercise) could quantify the economic consequences of selective attrition for prize money redistribution and labour-market signalling.

Professional cycling offers an unusually transparent window into economic decision-making under physical constraints. The findings suggest that even the decision to stop racing, which may appear purely physical, is shaped by the same forces of incentives, expected returns, and continuation costs that economic theory predicts.
